\section{Xử lý chuỗi}
\begingroup\small
\begin{multicols}{2}
\subsection{Z-function}
$Z[i]$ là độ dài tiền tố chung dài nhất của $s$ và
$s[i\ldots n-1]$; quy ước $Z[0]=n$.
Duy trì đoạn $[l,r)$ khớp với tiền tố. Nếu $i<r$,
lấy cận dưới $\min(r-i,Z[i-l])$, rồi so sánh thêm;
nếu vượt $r$, cập nhật đoạn. Mỗi lần so sánh mở rộng làm
$r$ tăng nên tổng thời gian $O(n)$.
Để tìm mẫu $p$ trong văn bản $t$, tính Z của
$p+\#+t$ với ký tự phân cách không có trong cả hai;
vị trí có Z ít nhất $|p|$ là một lần khớp. Nếu bảng chữ
không có ký tự phân cách, dùng chuỗi số nguyên với
giá trị sentinel riêng.

\subsection{Thuật toán Manacher}
Đặt $d_1[i]$ là số bán kính đối xứng lẻ quanh ký tự $i$:
độ dài dài nhất là $2d_1[i]-1$.
Đặt $d_2[i]$ là số bán kính đối xứng chẵn với tâm giữa
$i-1$ và $i$: độ dài dài nhất là $2d_2[i]$.
Duy trì palindrome phải nhất $[l,r]$, dùng tâm đối xứng
của $i$ để khởi tạo bán kính rồi mở rộng. Tổng thời gian
$O(n)$. Số xâu con đối xứng, tính cả vị trí trùng nội dung,
là $\sum_i(d_1[i]+d_2[i])$.
Mẫu code trả hai mảng bán kính này.

\subsection{Trie}
Trie lưu các tiền tố thành đỉnh. Với bảng chữ cố định,
mỗi cạnh gắn một ký tự; chèn một xâu dài $m$ mất $O(m)$.
Đếm số xâu đi qua mỗi đỉnh để trả số xâu có một tiền tố,
và đếm số xâu kết thúc để phân biệt tiền tố với từ hoàn chỉnh.
Mẫu dưới chỉ hỗ trợ chữ thường \texttt{a}--\texttt{z},
cho phép chèn trùng và không hỗ trợ xóa; số nút tối đa bằng
$1+$ tổng độ dài các xâu đã chèn.
\cppsnippet{code/strings/basic.hpp}

\subsection{Thuật toán Aho--Corasick}
Xây Trie của nhiều mẫu, rồi BFS tạo suffix link: link của
mỗi đỉnh trỏ tới hậu tố dài nhất vẫn là tiền tố của một mẫu.
Khi thiếu cạnh ký tự $c$, chuyển theo link; các cạnh thiếu
có thể điền sẵn để mỗi ký tự văn bản xử lý $O(1)$.
Cộng số mẫu kết thúc ở các đỉnh trên chuỗi link để đếm mọi
lần xuất hiện, kể cả mẫu trùng và các lần chồng nhau.
Với bảng chữ cỡ $\sigma$, thời gian và bộ nhớ xây dựng
$O(\sigma L)$ cho $L$ nút; quét văn bản $O(|t|)$.
Mẫu code đếm tổng số lần khớp, chỉ nhận chữ thường
\texttt{a}--\texttt{z}; phải thêm hết mẫu trước khi gọi
\texttt{build()} đúng một lần rồi mới truy vấn, và không thêm mẫu sau đó.
Nếu thêm mẫu rỗng, nó khớp $|t|+1$ lần, gồm cả trước ký tự đầu.

\subsection{Suffix array}
Mảng hậu tố $SA$ sắp các vị trí $i$ theo thứ tự từ điển của
$s[i\ldots n-1]$. Thuật toán doubling xếp theo cặp lớp
tương đương của hai nửa độ dài $2^k$; thêm sentinel nhỏ hơn
mọi ký tự để biến thành sắp vòng, rồi bỏ sentinel.
Dùng counting sort theo lớp cho $O(n\log n)$ thời gian
và $O(n)$ bộ nhớ. Mẫu code nhận byte khác \texttt{0},
trả các vị trí đánh số từ $0$; xâu rỗng trả mảng rỗng.
LCP có thể tính thêm bằng Kasai nếu bài cần so sánh nhiều
hậu tố hoặc đếm xâu con phân biệt, nhưng không nên đưa
cả bộ nếu bài chỉ cần thứ tự hậu tố.

\paragraph{Mẫu Aho--Corasick và suffix array (C++17).}
\cppsnippet{code/strings/advanced.hpp}
\end{multicols}
\endgroup
